\section{Resultados Preliminares}
\label{sec:resultados}

O baseline canônico tem aproximadamente 3,5 milhões de parâmetros e foi
treinado em Wikipedia EN, com semente fixa e modo determinístico. A
Tabela~\ref{tab:linguagem} resume os resultados de linguagem registrados no
model card.

\begin{table}[h]
\centering
\begin{tabular}{lrrrr}
\toprule
Métrica & full & no\_gravity & no\_gamma & no\_variable\_dim \\
\midrule
Val Loss & 6.415 & 6.385 & 6.321 & 6.428 \\
Val PPL & 611.2 & 593.2 & 556.1 & 619.2 \\
Train Loss & 6.344 & 5.883 & 6.252 & 6.308 \\
Tokens/s & 86163 & 92464 & 86468 & 86360 \\
\bottomrule
\end{tabular}
\caption{Resultados preliminares do baseline small\_3.5M e ablações.}
\label{tab:linguagem}
\end{table}

Em escala tão pequena, a diferença entre variantes ainda é marginal. A
variante \texttt{no\_gamma} obteve a menor perplexidade de validação em algumas
rodadas, sugerindo que o gamma-scaling pode exigir mais dados ou escala para se
tornar benéfico em modelagem de linguagem.

A avaliação topológica atual usa 99.840 vetores DRM, $K=80$ células de
Voronoi, 1.200 pontos para homologia, 10 restarts, razão de barra longa
$\rho=0{,}80$, projeção PCA3 e remoção dos 10\% pontos mais esparsos antes da
homologia persistente. O critério estrito de validação toroidal é:
\begin{equation}
  H_1=2,\qquad H_2=1,\qquad f_{T^2}\geq 0{,}60.
\end{equation}

Esse critério estrito ainda não foi atingido. Entretanto, os controles sugerem
que o treinamento desloca o manifold DRM de uma topologia inicial de gênero
alto e instável para uma assinatura homológica compatível com uma superfície
fechada quase-toroidal. A inicialização aleatória apresentou topologia de
gênero alto, $H_1=5$, $H_2=2$, com $f_{T^2}=0{,}00$ e \emph{Near T2}=0{,}00.
Após treinamento, três sementes independentes produziram \emph{Near T2}
elevado:
\begin{center}
\begin{tabular}{lrrrr}
\toprule
Controle & $H_1$ & $H_2$ & $f_{T^2}$ & Near T2 \\
\midrule
random init & 5 & 2 & 0{,}00 & 0{,}00 \\
seed 42 & 2 & 1 & 0{,}40 & 0{,}70 \\
seed 123 & 1 & 1 & 0{,}20 & 0{,}60 \\
seed 2025 & 2 & 1 & 0{,}30 & 0{,}80 \\
\bottomrule
\end{tabular}
\end{center}

Em duas de três sementes treinadas, a topologia mediana coincide com
$H_1=2$, $H_2=1$, assinatura esperada de $T^2$; em todas as sementes treinadas,
há sinal \emph{Near T2}, enquanto a inicialização aleatória não apresenta esse
sinal. Portanto, a leitura correta não é ``toro estável'', mas sim assinatura
homológica compatível com uma superfície fechada quase-toroidal e robusta a
sementes. O resultado também sugere que a assinatura não é mero artefato da
inicialização, pois o controle não treinado não exibiu sinal próximo de $T^2$.
